\scanpage{17}

formé des facettes incidentes à $s$) est un iso. En particulier $G_s\simeq\mathbb D_p$

b$'$) Soit $f\in F(\Pi)$. Alors $G_f\to\operatorname{Aut}(\Phi_{\leq f})$ (où $\Phi_{\leq f}$ est le polygône combinatoire à $q$ côtés des facettes incidentes à $f$) est un iso. En particulier, $G_f\simeq D_q$.

\unsure{c|b}) Soit $a\in A(\Pi)$. \deleted{Alors} Soit $\sigma$ \inserted{($\varphi$)} l'ens des sommets (faces) incidents à $a$. -- donc $\sigma,\varphi$ sont des ens. de card. 2.

Considérons

\[
 G_a\longrightarrow\unsure{\mathfrak S|\mathcal S}_{\sigma}\times\mathfrak S_{\varphi}=\mathbb Z/2\mathbb Z\times\mathbb Z/2\mathbb Z
\]

C'est un \uline{iso}.

\deleted{\uline{Exer} \uline{Proposition} Soit $\Pi$ polyèdre \inserted{comb.} régulier connexe fini, Pour tt $s\in S$,}

% archive sheet number: 16

